\slajdid{02-s026}
\begin{frame}[label=02-s026]{NILI kola su funkcionalno potpuna}
\setlength{\parskip}{0pt}
\setlength{\baselineskip}{12pt}
% element: 02-s026:e04
\alert{Svaku kombinacionu mrežu možemo realizovati samo NILI kolima.}
Invertor pri tome smatramo NILI kolom sa jednim ulazom.
\par
% element: 02-s026:e01
\Oznaka{Za realizaciju funkcije u obliku proizvoda zbirova}
% element: 02-s026:e02
\par
\noindent\(F=(C+B+A)(C+B+\bar A)(C+\bar B+A)(\bar C+B+A)(\bar C+\bar B+\bar A)\)\par
\noindent\(\phantom{F} =\overline{\overline{(C+B+A)(C+B+\bar A)(C+\bar B+A)(\bar C+B+A)(\bar C+\bar B+\bar A)}}\)\par
\noindent\(\phantom{F} =\overline{\overline{C+B+A}+\overline{C+B+\bar A}+\overline{C+\bar B+A}+\overline{\bar C+B+A}+\overline{\bar C+\bar B+\bar A}}\).\par
\vspace{.5mm}

% element: 02-s026:e03
\centering
\begin{tikzpicture}[font=\fontsize{9}{10.5}\selectfont,x=6mm,y=2.2mm,baseline=(current bounding box.center),
 DEinvertor/.style={not gate US,draw,scale=.6,minimum width=5mm,minimum height=4mm},
 DEprvinivo/.style={nor gate US,draw,scale=.6,logic gate inputs=nnn,rotate=-90,minimum width=5mm,minimum height=5mm},
 DEdrug nivo/.style={and gate US,draw,scale=.6,logic gate inputs=iiiii,rotate=-90,minimum width=6mm,minimum height=10mm}]
\def\DEkrajvodova{8.7}
\def\DEprviY{1.8}\def\DEdrugiY{-3.0}
% Gornji vod para je prava vrednost; donji je komplementna.
\ifdefined\DEparoviVodova\else
\def\DEparoviVodova{7.9/7.1/C,6.0/5.2/B,4.1/3.3/A}
\fi
\foreach \y/\low/\lab in \DEparoviVodova{
 \node[DEinvertor] (n\lab) at (.65,\y) {};
 \node[DEinvertor] (m\lab) at (2.0,\y) {};
 \draw[DEvod] (0,\y) node[left] {\(\lab\)}--(n\lab.input);
 \draw[DEvod] (n\lab.output)--(m\lab.input);
 \draw[DEvod] (m\lab.output)--(\DEkrajvodova,\y);
 \coordinate (t\lab) at ($(n\lab.output)!.5!(m\lab.input)$);
 \draw[DEvod] (t\lab)--(t\lab |- 0,\low)--(\DEkrajvodova,\low);
 \node[junction] at (t\lab) {};
}

\ifdefined\DEprviY\else\def\DEprviY{1.7}\fi
\ifdefined\DEdrugiY\else\def\DEdrugiY{-2.6}\fi
\foreach \i/\x in {1/3.1,2/4.4,3/5.7,4/7,5/8.3}{
 \node[DEprvinivo] (g\i) at (\x,\DEprviY) {};
}
\ifdefined\DEposUlazi\else
\def\DEposUlazi{1/3/7.9,1/2/6.0,1/1/4.1,2/3/7.9,2/2/6.0,2/1/3.3,3/3/7.9,3/2/5.2,3/1/4.1,4/3/7.1,4/2/6.0,4/1/4.1,5/3/7.1,5/2/5.2,5/1/3.3}
\fi
\foreach \i/\j/\y in \DEposUlazi{
 \draw[DEvod] (g\i.input \j)--(g\i.input \j |- 0,\y);
 \node[junction] at (g\i.input \j |- 0,\y) {};
}
\node[DEdrug nivo] (z) at (5.7,\DEdrugiY) {};
\coordinate (DEposUnutra) at ($(g3.output)!.33!(z.input 3)$);
\coordinate (DEposSpolja) at ($(g3.output)!.67!(z.input 3)$);
\draw[DEvod] (g1.output)--(g1.output |- DEposSpolja)-|(z.input 5);
\draw[DEvod] (g2.output)--(g2.output |- DEposUnutra)-|(z.input 4);
\draw[DEvod] (g3.output)--(z.input 3);
\draw[DEvod] (g4.output)--(g4.output |- DEposUnutra)-|(z.input 2);
\draw[DEvod] (g5.output)--(g5.output |- DEposSpolja)-|(z.input 1);
\draw[DEvod] (z.output)--++(0,-.3) node[below] {\(F\)};

\end{tikzpicture}

\(\longrightarrow\)
\begin{tikzpicture}[font=\fontsize{9}{10.5}\selectfont,x=6mm,y=2.2mm,baseline=(current bounding box.center),
 DEinvertor/.style={not gate US,draw,scale=.6,minimum width=5mm,minimum height=4mm},
 DEprvinivo/.style={nor gate US,draw,scale=.6,logic gate inputs=nnn,rotate=-90,minimum width=5mm,minimum height=5mm},
 DEdrug nivo/.style={nor gate US,draw,scale=.6,logic gate inputs=nnnnn,rotate=-90,minimum width=6mm,minimum height=10mm}]
\def\DEkrajvodova{8.7}
\def\DEprviY{1.8}\def\DEdrugiY{-3.0}
% Gornji vod para je prava vrednost; donji je komplementna.
\ifdefined\DEparoviVodova\else
\def\DEparoviVodova{7.9/7.1/C,6.0/5.2/B,4.1/3.3/A}
\fi
\foreach \y/\low/\lab in \DEparoviVodova{
 \node[DEinvertor] (n\lab) at (.65,\y) {};
 \node[DEinvertor] (m\lab) at (2.0,\y) {};
 \draw[DEvod] (0,\y) node[left] {\(\lab\)}--(n\lab.input);
 \draw[DEvod] (n\lab.output)--(m\lab.input);
 \draw[DEvod] (m\lab.output)--(\DEkrajvodova,\y);
 \coordinate (t\lab) at ($(n\lab.output)!.5!(m\lab.input)$);
 \draw[DEvod] (t\lab)--(t\lab |- 0,\low)--(\DEkrajvodova,\low);
 \node[junction] at (t\lab) {};
}

\ifdefined\DEprviY\else\def\DEprviY{1.7}\fi
\ifdefined\DEdrugiY\else\def\DEdrugiY{-2.6}\fi
\foreach \i/\x in {1/3.1,2/4.4,3/5.7,4/7,5/8.3}{
 \node[DEprvinivo] (g\i) at (\x,\DEprviY) {};
}
\ifdefined\DEposUlazi\else
\def\DEposUlazi{1/3/7.9,1/2/6.0,1/1/4.1,2/3/7.9,2/2/6.0,2/1/3.3,3/3/7.9,3/2/5.2,3/1/4.1,4/3/7.1,4/2/6.0,4/1/4.1,5/3/7.1,5/2/5.2,5/1/3.3}
\fi
\foreach \i/\j/\y in \DEposUlazi{
 \draw[DEvod] (g\i.input \j)--(g\i.input \j |- 0,\y);
 \node[junction] at (g\i.input \j |- 0,\y) {};
}
\node[DEdrug nivo] (z) at (5.7,\DEdrugiY) {};
\coordinate (DEposUnutra) at ($(g3.output)!.33!(z.input 3)$);
\coordinate (DEposSpolja) at ($(g3.output)!.67!(z.input 3)$);
\draw[DEvod] (g1.output)--(g1.output |- DEposSpolja)-|(z.input 5);
\draw[DEvod] (g2.output)--(g2.output |- DEposUnutra)-|(z.input 4);
\draw[DEvod] (g3.output)--(z.input 3);
\draw[DEvod] (g4.output)--(g4.output |- DEposUnutra)-|(z.input 2);
\draw[DEvod] (g5.output)--(g5.output |- DEposSpolja)-|(z.input 1);
\draw[DEvod] (z.output)--++(0,-.3) node[below] {\(F\)};

\end{tikzpicture}
\par
\end{frame}
